\part{Functorial Translation}

\section{Categories of Representational Systems}

A comparison of two contextual representation systems needs a notion of structure-preserving map, and the choice of which components to preserve, and in what sense, determines everything that follows. Let $\fR=(\cC,X,\Omega,\cA,\delta)$ with alphabet $\Sigma$ and $\fS=(\cC',X',\Omega',\cA',\delta')$ with alphabet $\Sigma'$ be systems, with observation maps $o$ and $o'$.

\begin{definition}[Translation]\label{def:translation}
A \emph{translation} $\Phi=(F,f,\eta,\alpha):\fR\to\fS$ consists of
\begin{enumerate}[label=(T\arabic*)]
\item a functor $F:\cC\to\cC'$;
\item a map $f:X\to X'$ of states;
\item maps $\eta_c:\Omega(c)\to\Omega'(Fc)$ for every context $c$, natural in the sense that $\eta_c\circ\Omega(u)=\Omega'(Fu)\circ\eta_d$ for every arrow $u:c\to d$;
\item a map $\alpha:\Sigma\to\Sigma'$ of actions,
\end{enumerate}
subject to the following conditions.
\begin{enumerate}[label=(\Alph*)]
\item[(O)] \emph{Observation preservation}: $o'_{Fc}\circ f=\eta_c\circ o_c$ for every context $c$.
\item[(A)] \emph{Lax authority preservation}: $\alpha(\cA(c))\subseteq\cA'(Fc)$ for every context $c$.
\item[(D)] \emph{Simulation of dynamics}: if $\delta(x,a)=y$ then $\delta'(f(x),\alpha(a))=f(y)$.
\end{enumerate}
A translation is \emph{observational} if only (T1) to (T3) and (O) are required. It is \emph{authority-reflecting} if $\cA(c)=\alpha^{-1}(\cA'(Fc))$, and \emph{strong} if in addition $\delta(x,a)$ is defined exactly when $\delta'(f(x),\alpha(a))$ is.
\end{definition}

The minimal requirement for the mapping on contexts to qualify as a functor is that it preserve identities and composition, which is not automatic for an arbitrary assignment of arrows. Where the category of contexts is generated by a graph of navigation edges, freeness gives a convenient way to meet this requirement, since a map of generating edges to arrows of the target extends uniquely to a functor on the path category. The conditions (A) and (D) are \emph{lax}: they require inclusion and implication in one direction, not equality, and the strong variants are the strict counterparts. The observation condition (O) is strict.

\begin{proposition}[The category $\CRS$; \stN]\label{prop:crs}
Translations compose by $(F',f',\eta',\alpha')\circ(F,f,\eta,\alpha)=(F'F,\,f'f,\,\eta'_{F(-)}\circ\eta_{(-)},\,\alpha'\alpha)$, with identities $(\id,\id,\id,\id)$, and this makes systems and translations into a category $\CRS$. The authority-reflecting and strong translations are closed under composition.
\end{proposition}
\begin{proof}
For naturality of the composite $\eta'_{Fc}\eta_c$, take $u:c\to d$ and compute $\eta'_{Fc}\eta_c\Omega(u)=\eta'_{Fc}\Omega'(Fu)\eta_d=\Omega''(F'Fu)\eta'_{Fd}\eta_d$, using naturality of $\eta$ and then of $\eta'$ at the arrow $Fu$. For (O), $o''_{F'Fc}f'f=\eta'_{Fc}o'_{Fc}f=\eta'_{Fc}\eta_co_c$. For (A), $\alpha'\alpha(\cA(c))\subseteq\alpha'(\cA'(Fc))\subseteq\cA''(F'Fc)$. For (D), apply the hypothesis twice. Identities satisfy all conditions, and associativity is inherited from composition of maps and functors. For the authority-reflecting case, $(\alpha'\alpha)^{-1}(\cA''(F'Fc))=\alpha^{-1}(\alpha'^{-1}(\cA''(F'Fc)))=\alpha^{-1}(\cA'(Fc))=\cA(c)$, and the strong case is the same argument applied to definedness.
\end{proof}

The first preservation result is that every translation preserves observational indistinguishability, and the converse needs more.

\begin{theorem}[Translation preserves source indistinguishability; \stN]\label{thm:presindist}
For every translation $\Phi:\fR\to\fS$, every context $c$ of $\fR$, and all states $x,y$ of $\fR$,
\[
x\sk_c y\;\Longrightarrow\;f(x)\sk'_{Fc}f(y).
\]
\end{theorem}
\begin{proof}
If $o_c(x)=o_c(y)$ then $o'_{Fc}f(x)=\eta_co_c(x)=\eta_co_c(y)=o'_{Fc}f(y)$ by (O).
\end{proof}

\begin{proposition}[Reflection requires additional structure; \stN]\label{prop:reflection}
The converse of Theorem~\ref{thm:presindist} does not hold in general. It holds at $c$ for every translation in which $\eta_c$ is injective on $o_c(X)$, and conversely if the converse holds at $c$ then $\eta_c$ is injective on $o_c(X)$.
\end{proposition}
\begin{proof}
For the failure take $X=X'=\{0,1\}$, $f=\id$, $\Omega(c)=\{0,1\}$ with $o_c=\id$, and $\Omega'(c)=\{\ast\}$ with $\eta_c$ constant. Then $f(0)\sk'_cf(1)$ while $0\not\sk_c1$. If $\eta_c$ is injective on $o_c(X)$ and $f(x)\sk'_{Fc}f(y)$, then $\eta_co_c(x)=o'_{Fc}f(x)=o'_{Fc}f(y)=\eta_co_c(y)$, so $o_c(x)=o_c(y)$. Conversely, if the converse holds and $\eta_c(o_c(x))=\eta_c(o_c(y))$, then $f(x)\sk'_{Fc}f(y)$ by (O), hence $x\sk_cy$ and $o_c(x)=o_c(y)$.
\end{proof}

Preservation and reflection of distinctions are thus mathematically separate. A translation that merges the observation space of a context collapses distinctions, and no property of the target can recover them.

\section{Common Abstractions and Forgetful Functors}

Two systems have a \emph{common abstraction} when both translate into a third system. Forgetting components yields natural families of such abstractions. The reduct that forgets $\cA$ and $\delta$ sends a system to its observational structure $(\cC,X,\Omega,o)$; the reduct that forgets $\delta$ alone retains authority; and each is a functor from $\CRS$, since the conditions defining translation restrict to the retained components. A translation can be classified by what it covers and what it retains.

\begin{definition}[Kinds of representation]\label{def:kinds}
A translation $\Phi=(F,f,\eta,\alpha)$ is \emph{state-surjective} if $f$ is surjective, \emph{context-surjective} if $F$ is surjective on objects, and \emph{observation-reflecting} if $\eta_c$ is injective on $o_c(X)$ for every $c$, equivalently by Proposition~\ref{prop:reflection} if $x\sk_cy\Leftrightarrow f(x)\sk'_{Fc}f(y)$ for all $c$. It is a \emph{surjective representation} if it is state-surjective and context-surjective, a \emph{faithful representation} if it is moreover observation-reflecting, and an \emph{isomorphism of systems} if it has an inverse translation in $\CRS$.
\end{definition}

The first result states what a common surjective abstraction does and does not yield. It is valid and elementary, and its modesty is the point: an abstraction retains only the distinctions that its observation maps carry.

\begin{theorem}[Common abstraction without equivalence; \stN]\label{thm:central}
Let $\fR_1$ and $\fR_2$ be contextual representation systems admitting observational translations $\Phi_i=(F_i,f_i,\eta^i):\fR_i\to\fS$ into a common system $\fS$, and suppose the state maps are surjective, $f_i:X_i\twoheadrightarrow X_{\fS}$. Then:
\begin{enumerate}[label=(\roman*)]
\item every state of $\fS$ has at least one representative in each $X_i$;
\item for every context $c$, $x\sk^i_cy$ implies $f_i(x)\sk^{\fS}_{F_ic}f_i(y)$, and target observations pull back consistently: if $s\in X_{\fS}$ is represented by $x_1\in X_1$ and $x_2\in X_2$, and $F_1c_1=F_2c_2=e$, then $\eta^1_{c_1}o^1_{c_1}(x_1)=\eta^2_{c_2}o^2_{c_2}(x_2)=o^{\fS}_e(s)$;
\item $\fR_1$ and $\fR_2$ need not be isomorphic, need not have state sets of the same cardinality, and need not have isomorphic context categories.
\end{enumerate}
\end{theorem}
\begin{proof}
Part (i) is surjectivity and part (ii) is Theorem~\ref{thm:presindist} together with condition (O). For (iii), let $\fS$ be the system with one context, one state, a trivial observation space, and no actions. Let $\fR_1$ have one state and $\fR_2$ two states, each with one context, trivial observations, and no actions, so that both map onto $\fS$ by constant maps, yet $|X_1|\neq|X_2|$ rules out an isomorphism. For context categories, let $\cC_1$ be the terminal category and $\cC_2$ the one-object category whose endomorphisms are $\{1,e\}$ with $e^2=e$, each with a one-point state set, trivial observations, and no actions. Both admit functors to the terminal category, and $\cC_1\not\cong\cC_2$.
\end{proof}

Common abstraction is cheap in both directions. There is a terminal system, so that every system translates to it and every pair of systems has a common abstraction that retains nothing, and there is a disjoint-union system that retains everything and identifies nothing.

\begin{proposition}[Trivial common abstractions; \stN]\label{prop:trivialabs}
$\CRS$ has a terminal object $\fT$, with one context, one state, a one-point observation space, one action $\sigma$ admissible at the context, and $\delta(\ast,\sigma)=\ast$. Consequently any two systems have a common state-surjective and context-surjective abstraction.
\end{proposition}
\begin{proof}
For any $\fR$ there is exactly one functor $\cC\to1$, one map $X\to\{\ast\}$, one family of maps $\Omega(c)\to\{\ast\}$, and one map $\Sigma\to\{\sigma\}$. Condition (O) holds because the target has one point, (A) holds because $\{\sigma\}\subseteq\cA'(\ast)$, and (D) holds because $\delta'(\ast,\sigma)=\ast=f(y)$. Uniqueness of each component gives uniqueness of the translation.
\end{proof}

A significant common abstraction therefore has to be singled out by what it retains. If a family of observations is fixed, there is a canonical choice.

\begin{theorem}[Terminal observed system; \stE, applied here]\label{thm:terminalobs}
Fix a small category $\cC$ and a presheaf $\Omega:\cC^{\op}\to\Set$. An \emph{$\Omega$-observed set} is a pair $(X,o)$ with $o:\Delta X\Rightarrow\Omega$ natural, and a morphism $h:(X,o)\to(X',o')$ is a map with $o'h=o$. Let $\Gamma(\Omega)=\lim\Omega$ be the set of global sections, that is, families $s=(s_c)_c\in\prod_c\Omega(c)$ with $\Omega(u)(s_d)=s_c$ for every arrow $u:c\to d$, and let $\mathrm{ev}_c:\Gamma(\Omega)\to\Omega(c)$ be the projections. Then $(\Gamma(\Omega),\mathrm{ev})$ is the terminal $\Omega$-observed set, and the unique morphism from $(X,o)$ is the \emph{profile map}
\[
\pi:X\to\Gamma(\Omega),\qquad\pi(x)=(o_c(x))_{c\in\Ob(\cC)}.
\]
\end{theorem}
\begin{proof}
By coherence $\Omega(u)(o_d(x))=o_c(x)$ for $u:c\to d$, so $\pi(x)$ is a global section, and $\mathrm{ev}_c\pi=o_c$. If $h$ is any morphism, $\mathrm{ev}_ch(x)=o_c(x)$ for every $c$, which determines $h(x)$ as the family $\pi(x)$.
\end{proof}

\begin{corollary}[What the canonical abstraction retains and forgets; \stN]\label{cor:retains}
For an $\Omega$-observed set with profile map $\pi$:
\begin{enumerate}[label=(\alph*)]
\item $\pi$ is injective exactly when the combined observations separate states, and $X/\bigcap_c{\sk_c}\cong\im\pi$, which identifies the global indistinguishability quotient of Proposition~\ref{prop:globalindist} with the set of realized profiles;
\item $\pi$ is surjective exactly when every global section of $\Omega$ is realized by a state;
\item two $\Omega$-observed sets $(X_1,o^1)$ and $(X_2,o^2)$ have the same realized profiles, $\im\pi_1=\im\pi_2$, exactly when they have the same canonical abstraction, and they are isomorphic as $\Omega$-observed sets exactly when $|\pi_1^{-1}(s)|=|\pi_2^{-1}(s)|$ for every $s\in\Gamma(\Omega)$.
\end{enumerate}
\end{corollary}
\begin{proof}
Parts (a) and (b) restate Proposition~\ref{prop:globalindist} and surjectivity. For (c), an isomorphism $h$ satisfies $\pi_2h=\pi_1$, so it carries each fiber of $\pi_1$ bijectively to the fiber of $\pi_2$ over the same profile, and conversely fiberwise bijections assemble to an isomorphism.
\end{proof}

The corollary identifies precisely what the canonical abstraction forgets about the observed set, namely the multiplicity of microstates over each realized profile, and what remains to be compared separately, namely $\cA$ and $\delta$. Part (b) is the global form of the existence half of the gluing condition and part (a) the global form of separation. The difference between a common abstraction and an equivalence is made concrete by the following pair of systems.

\begin{example}[Same observational abstraction, no translation back; \stN]\label{ex:flip}
Let $\cC=\{c\to d\}$ with $\Omega(d)=\{0,1\}^2$, $\Omega(c)=\{0,1\}$, and $\Omega(c\to d)$ the first projection. Let $\Sigma=\{\texttt{flip}\}$, $\cA(c)=\varnothing$, and $\cA(d)=\{\texttt{flip}\}$. Let $\fR_1$ have $X_1=\{0,1\}^2$, $o_d=\id$, $o_c$ the first coordinate, and $\delta_1(x,\texttt{flip})=(1-x_1,x_2)$. Let $\fR_2$ have $X_2=\{0,1\}^2\times\{0,1,2\}$ with $o_d(x,k)=x$, $o_c(x,k)=x_1$, and $\delta_2((x,k),\texttt{flip})=((1-x_1,x_2),k+1)$ when $k<2$ and undefined when $k=2$. The projection $\fR_2\to\fR_1$ forgetting the counter, with $F=\id$, $\eta=\id$, and $\alpha=\id$, is a translation that is state-surjective and observation-reflecting, and both systems realize every profile, so their canonical abstractions coincide. There is, however, no translation $\fR_1\to\fR_2$ at all. If $h$ were the state map, the run $x_0\to x_1\to x_0$ of $\fR_1$ with $x_0=(0,0)$ and $x_1=(1,0)$ would be sent by (D) to a run $h(x_0)\to h(x_1)\to h(x_0)$ of $\fR_2$, but $\delta_2$ increases the counter at every step, which is impossible. Hence $\fR_1$ and $\fR_2$ are not isomorphic in $\CRS$.
\end{example}

The example separates the observational abstraction from the dynamics: $\fR_2$ can be flipped at most twice, and the abstraction does not see this. This is the sense in which an abstraction established relative to observations can hold while the systems differ in operational behavior.

A second family of comparisons connects the structural examples of Part~II. Let $T$ be a finite rooted tree, ordered by nesting with the root greatest, and let $L=(L(p))_{p\in T}$ be a family of local sets. The pair $(T,L)$ is a \emph{decorated rooted tree}, and it generates two systems on the common state set $X=\prod_pL(p)$, with $\cA=\varnothing$ and no transitions. The \emph{region system} $\fR^{\downarrow}(T,L)$ has $\cC^{\downarrow}=(T,\preceq)$ with an arrow $q\to p$ when $q\preceq p$, observation space $\Omega^{\downarrow}(p)=\prod_{q\preceq p}L(q)$ over the nodes nested in $p$, and $o_p$ the projection. This is the filesystem viewed through subtrees, and a world viewed through its regions. The \emph{accumulation system} $\fR^{\uparrow}(T,L)$ has $\cC^{\uparrow}=(T,\preceq)^{\op}$, observation space $\Omega^{\uparrow}(q)=\prod_{p\succeq q}L(p)$ over the ancestors of $q$, and $o_q$ the projection, as in Example~\ref{ex:accum}. This is lexical scope. Both are coherent.

\begin{proposition}[Orientation duality; \stN]\label{prop:orientation}
Let $(T,L)$ be a decorated rooted tree with all $L(p)$ nonempty.
\begin{enumerate}[label=(\alph*)]
\item The nesting categories of the filesystem model and of the lexical-scope model, each generated by the parent relation, are isomorphic whenever the parent trees are isomorphic.
\item The profile maps of $\fR^{\downarrow}(T,L)$ and $\fR^{\uparrow}(T,L)$ are both bijections onto their global sections, so both systems have the identity as global indistinguishability relation, and $\Gamma(\Omega^{\downarrow})\cong X\cong\Gamma(\Omega^{\uparrow})$.
\item The refinement categories $\cC^{\downarrow}=(T,\preceq)$ and $\cC^{\uparrow}=(T,\preceq)^{\op}$ are isomorphic if and only if $T$ is a chain. In particular, if $T$ has a node with two children, no isomorphism of context categories carries the region system to the accumulation system.
\end{enumerate}
\end{proposition}
\begin{proof}
(a) is the statement that the parent relation has the same orientation in both models. For (b), a global section of $\Omega^{\downarrow}$ is determined by its component at the root, which lies in $\prod_qL(q)=X$ and restricts to every other component, so $\Gamma(\Omega^{\downarrow})\cong X$ and $\pi$ is the identity under this identification. For $\Omega^{\uparrow}$, every node is an ancestor of some leaf, and a compatible family of ancestor chains over all leaves agrees on shared ancestors, hence defines a unique element of $\prod_pL(p)$, so $\Gamma(\Omega^{\uparrow})\cong X$ and $\pi$ is again the identity. For (c), the poset $(T,\preceq)$ has exactly one maximal element, the root, and as many minimal elements as leaves. A chain is isomorphic to its dual. If $T$ has a node with two children then it has at least two leaves, so its dual has at least two maximal elements and one minimal element, and an order isomorphism must preserve the number of maximal elements.
\end{proof}

This clarifies a tempting formulation: a filesystem and a lexical scope structure are the same tree. They share the parent tree and the state space, and they share their global observational abstraction, but the information flows in opposite directions along it. The two systems differ in the lattice of local partitions: the region system has $\sk^{\downarrow}_p$ identifying states that agree on the subtree of $p$, and the accumulation system has $\sk^{\uparrow}_q$ identifying states that agree on the ancestor chain of $q$, and these two families are not related by any isomorphism of the indexing categories when $T$ branches. The common abstraction is the decorated tree with its undirected rooted structure, and the orientation of information flow is an additional structure that the abstraction does not share. Nor does a shared decorated tree establish equivalence of file operations and name-resolution semantics, which differ in $\cA$ and $\delta$. Moreover, a filesystem and a game world arise from the same decorated tree. A hub with two rooms, each carrying one local bit, is the region system of a directory with two subdirectories each carrying one file attribute, and the two have identical observational structure and different dynamics.

\begin{center}
\begin{tabular}{lll}
\toprule
 & Filesystem on $(T,L)$ & Game world on $(T,L)$\\
\midrule
contexts & directories & hub and rooms\\
local content $L(p)$ & entry attributes & room object states\\
observation at $p$ & listing of subtree & view of region\\
actions $\Sigma$ & \texttt{cd}, \texttt{chmod}, \texttt{rm} & move, take, toggle\\
authority $\cA$ & mode bits and ownership & key and door conditions\\
\bottomrule
\end{tabular}
\end{center}

\section{Surjectivity, Fullness, Faithfulness, and Equivalence}

The vocabulary of functors separates several properties that informal talk about mappings between systems runs together. A functor $F:\cC\to\cD$ is \emph{full} if it is surjective on each hom-set, \emph{faithful} if it is injective on each hom-set, and \emph{essentially surjective} if every object of $\cD$ is isomorphic to some $F(c)$. Surjectivity on objects is stronger than essential surjectivity and independent of the other two.

\begin{theorem}[Equivalence criterion; \stE]\label{thm:equiv}
A functor $F:\cC\to\cD$ is an equivalence of categories if and only if it is full, faithful, and essentially surjective \cite{maclane}.
\end{theorem}

This is the standard benchmark for the stronger equivalence claim. Everything weaker is a statement about a subset of these properties, and the next result shows that the weakest, surjectivity on objects, implies none of the others.

\begin{proposition}[Surjectivity does not imply equivalence; \stN]\label{prop:surjnotequiv}
A functor that is surjective on objects need not be full or faithful, and need not be an equivalence.
\end{proposition}
\begin{proof}
Let $\cD_2$ be the discrete category with two objects $a,b$ and let $1$ be the terminal category. The unique functor $\cD_2\to1$ is surjective on objects but not full, since $\Hom(a,b)=\varnothing$ while $\Hom(F a,Fb)=\{\id\}$. It is not an equivalence, because equivalent categories have the same number of isomorphism classes of objects. For faithfulness, let $M$ be the one-object category with endomorphisms $\{1,e\}$ with $e^2=e$. The functor $M\to1$ is surjective on objects and full, and not faithful because $1$ and $e$ have the same image. It is not an equivalence because corresponding hom-sets of equivalent categories have the same cardinality.
\end{proof}

\begin{proposition}[No unrestricted faithful collapse; \stN]\label{prop:nocollapse}
A functor that identifies two distinct parallel arrows is not faithful. In particular the reflection $\Path(G)\to\Reach(G)$ of Proposition~\ref{prop:catfromctx} is not faithful whenever $G$ has two distinct routes between the same vertices, so no faithful translation of contexts can collapse a distinction between routes.
\end{proposition}
\begin{proof}
Faithfulness is injectivity on every hom-set, so two distinct arrows with the same source and target and the same image violate it.
\end{proof}

Rigidity of posets places a sharp limit on what categorical equivalence can establish for context hierarchies.

\begin{proposition}[Equivalence of posets is isomorphism; \stN]\label{prop:posetrigid}
A functor between posets, regarded as categories, is an equivalence if and only if it is an order isomorphism.
\end{proposition}
\begin{proof}
An order isomorphism is an equivalence. Conversely, an equivalence is full and faithful, so it reflects the order, $F(p)\le F(q)\Rightarrow p\le q$. If $F(p)=F(q)$ then $p\le q$ and $q\le p$, so $p=q$ by antisymmetry and $F$ is injective. In a poset the only isomorphisms are identities, so essential surjectivity is surjectivity on objects, and $F$ is a bijective order-reflecting monotone map.
\end{proof}

\begin{corollary}[Morita rigidity for posets; \stE]\label{cor:morita}
Two posets have equivalent categories of presheaves if and only if they are isomorphic.
\end{corollary}
\begin{proof}
Two small categories have equivalent presheaf categories exactly when their Cauchy completions are equivalent \cite{borceux}. A poset has no idempotents other than identities, so its Cauchy completion is equivalent to itself, and Proposition~\ref{prop:posetrigid} gives isomorphism.
\end{proof}

Hence nontrivial identification of context hierarchies cannot come from categorical equivalence of posets, and the presheaf semantics of Part~III does not blur the distinction between different hierarchies. Nontrivial identification must come from functors that coarsen, which are neither full nor faithful, or from relations weaker than equivalence, such as simulation, bisimulation, and observational equivalence. This explains the shape of the results: they are statements of common abstraction together with an analysis of what the abstraction preserves.

The most useful positive result concerns what a state-surjective, observation-reflecting translation does establish.

\begin{theorem}[Relative equivalence; \stN]\label{thm:relequiv}
Let $\Phi=(F,f,\eta,\alpha):\fR\to\fS$ be a translation that is state-surjective and observation-reflecting. Then for every context $c$ the map
\[
f_c:Q(c)\to Q'(Fc),\qquad[x]_c\mapsto[f(x)]_{Fc}
\]
is a bijection, and the family $(f_c)$ is a natural isomorphism $Q\cong F^{\ast}Q'$ of presheaves on $\cC$, where $F^{\ast}Q'=Q'\circ F^{\op}$.
\end{theorem}
\begin{proof}
The map is well defined by Theorem~\ref{thm:presindist} and injective because $\Phi$ is observation-reflecting. It is surjective because $f$ is. For $u:c\to d$, $f_cQ(u)[x]_d=[f(x)]_{Fc}=Q'(Fu)f_d[x]_d$.
\end{proof}

\begin{corollary}[Equivalence relative to retained structure; \stN]\label{cor:relequiv}
If $\fR_1$ and $\fR_2$ admit state-surjective, observation-reflecting translations $\Phi_i:\fR_i\to\fS$ with $F_i=\id$ on a common $\cC$, then $Q_{\fR_1}\cong Q_{\fR_2}$ as presheaves on $\cC$. The systems are then equivalent as observational presheaves, and may still differ in $X$, $\cA$, and $\delta$, as in Example~\ref{ex:flip}.
\end{corollary}

This is the form in which a common abstraction can establish equivalence without full equivalence: equivalence of the canonical presheaves, relative to the distinctions that the abstraction retains. The next proposition shows that sharing a common quotient does not even yield mutual translation of the systems.

\begin{proposition}[Common quotient does not imply mutual translation; \stN]\label{prop:commonquot}
\begin{enumerate}[label=(\alph*)]
\item There exist systems $\fR_1,\fR_2$ with a common state-surjective, observation-reflecting abstraction $\fS$ such that no translation $\fR_1\to\fR_2$ is bijective on states and both preserves and reflects observations.
\item There exist such systems with no translation $\fR_1\to\fR_2$ at all.
\end{enumerate}
\end{proposition}
\begin{proof}
(a) Let $\cC=1$ and $\Omega(\ast)=\{0,1\}$. Let $X_1=X_2=\{1,2,3,4\}$ with $o^1$ constant $0$ on $\{1,2\}$ and $1$ on $\{3,4\}$, and $o^2$ constant $0$ on $\{1\}$ and $1$ on $\{2,3,4\}$. Both map by $o^i$ onto $\fS$ with $X'=\{0,1\}$, $o'=\id$, and are surjective and observation-reflecting. A bijection $h:X_1\to X_2$ that both preserves and reflects observations would carry the observation classes $\{1,2\},\{3,4\}$ onto the classes $\{1\},\{2,3,4\}$, which have different sizes. (b) is Example~\ref{ex:flip}.
\end{proof}

Without the bijectivity requirement in (a), mutual translations exist: send the class $\{1,2\}$ to $\{1\}$ and $\{3,4\}$ into $\{2,3,4\}$, and conversely, each preserving and reflecting observations. The requirement is therefore needed to keep constant-like maps from trivializing the example, and part (b) shows that the dynamics supply a genuine obstruction independent of it.

\section{Translation Between Mathematical Languages}

Set-theoretic, type-theoretic, logical, algebraic, and categorical presentations of mathematical structure are connected by translations of different strengths, and the discussion above supplies a vocabulary for keeping them apart. A \emph{syntactic translation} from a theory $T$ to a theory $T'$ is a compositional map $\varphi\mapsto\varphi^{\ast}$ of formulas such that $T\vdash\varphi$ implies $T'\vdash\varphi^{\ast}$, and it \emph{reflects} provability when the converse holds. A \emph{semantic interpretation} is a functor between categories of models, typically a reduct. Two theories are \emph{definitionally equivalent} when they have a common definitional extension, and \emph{Morita equivalent} when, after adding new sorts for products, subsorts, and quotients, they are definitionally equivalent \cite{barretthalvorson}. Finally the categories of models may be \emph{equivalent as categories}. These notions are logically independent: a syntactic translation need not reflect provability, a semantic interpretation need not be an equivalence, and categorical equivalence of models need not come from any translation of syntax with preserved proof theory. The following examples show each separation concretely.

\begin{proposition}[Classical and intuitionistic arithmetic; \stE]\label{prop:PAHA}
Let $\mathrm{PA}$ be Peano arithmetic and $\mathrm{HA}$ Heyting arithmetic.
\begin{enumerate}[label=(\alph*)]
\item The G\"odel--Gentzen negative translation sends each sentence $\varphi$ to a sentence $\varphi^N$ such that $\mathrm{PA}\vdash\varphi$ implies $\mathrm{HA}\vdash\varphi^N$ \cite{troelstravandalen}.
\item $\mathrm{PA}$ is conservative over $\mathrm{HA}$ for $\Pi^0_2$ sentences \cite{friedman}.
\item $\mathrm{HA}$ does not prove every theorem of $\mathrm{PA}$: for a G\"odel sentence $G$ of $\mathrm{PA}$, $\mathrm{PA}\vdash G\vee\neg G$ while $\mathrm{HA}\nvdash G\vee\neg G$, assuming $\mathrm{PA}$ is consistent.
\end{enumerate}
\end{proposition}
\begin{proof}[Proof of (c)]
Excluded middle gives $\mathrm{PA}\vdash G\vee\neg G$. $\mathrm{HA}$ has the disjunction property \cite{troelstravandalen}: if $\mathrm{HA}\vdash A\vee B$ then $\mathrm{HA}\vdash A$ or $\mathrm{HA}\vdash B$. If $\mathrm{HA}\vdash G\vee\neg G$ then, since $\mathrm{PA}$ extends $\mathrm{HA}$, $\mathrm{PA}\vdash G$ or $\mathrm{PA}\vdash\neg G$, contradicting the independence of $G$.
\end{proof}

The same phenomenon appears algebraically. Classical propositional logic corresponds to Boolean algebras and intuitionistic propositional logic to Heyting algebras, and the double-negation map is the algebraic form of the negative translation.

\begin{proposition}[Glivenko map; \stE]\label{prop:glivenko}
For a Heyting algebra $H$, the map $a\mapsto\neg\neg a$ is a surjective Heyting algebra homomorphism $H\to H_{\neg\neg}$ onto the Boolean algebra of regular elements, and it is injective if and only if $H$ is Boolean \cite{glivenko}.
\end{proposition}

This is a surjective homomorphism of structures that is not an equivalence unless $H$ is already Boolean, an instance of Proposition~\ref{prop:surjnotequiv}: the translation retains the classical theory and discards the intuitionistic distinctions among elements with the same double negation.

There are also genuine equivalences whose presentations differ. Finitary algebraic theories presented as Lawvere theories correspond to finitary monads on $\Set$, with equivalent categories of models, although one presentation has operations and equations and the other has none \cite{lawvere}. Typed lambda calculi with products and exponentials, taken up to their equational theories, correspond to cartesian closed categories in an equivalence in which terms are morphisms \cite{lambekscott}; this equivalence does not see evaluation strategy or reduction sequences, since it is stated at the level of equality of terms. Rings $R$ and $S$ are Morita equivalent when their module categories are equivalent: a field $k$ and the matrix ring $M_2(k)$ are Morita equivalent and not isomorphic, and the equivalence preserves module-theoretic invariants such as the number of simple modules while losing ring-theoretic invariants such as commutativity. Constructive set theory is interpretable in Martin-L\"of type theory \cite{aczel}, and the strength of the type theory underlying contemporary proof assistants has been compared with set theory augmented by inaccessible cardinals \cite{carneiro}; these are results of the form ``interpretable in, with a stated strength'' and not statements of equivalence. In each case the relation among languages is specified by what the translation preserves.

The correspondence of Part~III has a syntactic counterpart in which the orientation duality of Proposition~\ref{prop:orientation} reappears. For a model $M$ and a set $V$ of variables, let $X=M^V$ be the set of valuations and let contexts be finite subsets $\Gamma\subseteq V$ with $\Omega(\Gamma)=M^{\Gamma}$ the restriction of valuations. This is a region system in the sense of Example~\ref{ex:region}. Terms in context $\Gamma$ form a family $\mathrm{Term}(\Gamma)$ with weakening maps $\mathrm{Term}(\Gamma)\to\mathrm{Term}(\Delta)$ for $\Gamma\subseteq\Delta$, which run in the direction opposite to restriction, and the evaluation of a weakened term on a valuation equals the evaluation of the original on its restriction. Syntax is covariant and semantics is contravariant, and the pairing between them is natural. The orientation of information flow is therefore a feature that distinguishes a presentation by terms from a presentation by valuations, as it distinguishes a lexical scope from a filesystem.

\begin{question}[Interpretations and preservation]\label{q:interp}
To what extent do properties of an interpretation between first-order theories, such as conservativity and bi-interpretability, correspond to properties of the induced translations of valuation systems, such as being observation-reflecting or state-surjective in the sense of Definition~\ref{def:kinds}?
\end{question}

\section{Universality and Its Limits}

Computational universality is a statement that a formalism computes every computable function. The strongest form of the observation that universal systems are interchangeable is Rogers' isomorphism theorem: any two acceptable numberings of the partial computable functions, which includes the programming systems of any two universal languages, are related by a computable bijection of program indices that preserves the function computed \cite{rogers}. Universality therefore yields a translation between any two universal systems, and the translation preserves exactly one thing, the input-output function of each program. In the terminology above, it is a surjective representation onto a one-observation system, the weakly terminal abstraction of Proposition~\ref{prop:trivialabs} specialized to the observation ``what function is computed''. Nothing else need be preserved, and in general nothing else is.

Several distinct kinds of structure are lost. Resource behavior depends on the model: translating between Turing machines and random-access machines costs a polynomial overhead, so complexity classes defined up to polynomial factors are stable, while fine-grained time bounds are not, and whether the stability extends to every physically realizable model is open, with quantum computation the principal candidate for violating it if the relevant complexity classes differ. Expressive power in a language-design sense is finer than computability: one language is more expressive than another when a construct cannot be eliminated by local rewriting without global reorganization of programs \cite{felleisen}. Definability at higher types separates languages that compute the same first-order functions: the parallel-or operator is not definable in the language PCF, though it is computable, and adding it yields the language for which the standard model is fully abstract \cite{plotkin}. Type systems and proof theories differ for languages with the same computable functions, as the contrast between simply typed and untyped lambda calculi illustrates.

Contextual equivalence provides the bridge to Part~III. Let $X$ be a set of program fragments, let the family $\cU$ of closing program contexts be a discrete set of observation contexts, and let the observation of a fragment $M$ at a program context $C[\cdot]$ be whether $C[M]$ terminates, valued in $\{0,1\}$ \cite{morris}. Then $\sk_\cU$ is contextual equivalence. A denotational valuation $[\![\cdot]\!]:X\to D$ is sound when $[\![M]\!]=[\![N]\!]$ implies $M\sk_\cU N$, which says that its image-quotient is $\cU$-sound in the sense of Section~11, and it is fully abstract when the converse also holds.

\begin{proposition}[Full abstraction is the canonical quotient; \stE]\label{prop:fullabs}
The quotient $X\to X/{\sk_\cU}$ by contextual equivalence factors uniquely through every sound surjective valuation $X\to Z$, as $X\to Z\to X/{\sk_\cU}$. A sound valuation is fully abstract exactly when its kernel is $\sk_\cU$, that is, when it is itself the terminal $\cU$-sound quotient of Theorem~\ref{thm:coarsest}.
\end{proposition}
\begin{proof}
This is Theorem~\ref{thm:coarsest}(b) applied to the discrete family $\cU$, together with the observation that a sound surjective valuation has kernel contained in $\sk_\cU$, with equality exactly in the fully abstract case.
\end{proof}

The proposition does not claim novelty for the concept of full abstraction, which is classical \cite{plotkin}. It records that full abstraction is an instance of the canonical quotient of the preceding sections, and that the universal property of Theorem~\ref{thm:coarsest} is a statement about every sound semantics at once.

\begin{proposition}[Universality does not fix the observation; \stE]\label{prop:univinsufficient}
Let $\mathrm{PCF}$ and $\mathrm{PCF}+\mathrm{por}$ denote the typed functional language with the usual arithmetic and fixed-point constructs and the same language extended by parallel-or. Both languages are universal for the first-order computable functions on the natural numbers, yet their contextual equivalences on higher-type terms differ, so the two canonical quotients of Theorem~\ref{thm:coarsest} are distinct.
\end{proposition}
\begin{proof}
The parallel-or operator is not definable in $\mathrm{PCF}$ \cite{plotkin}, and terms of higher type that no $\mathrm{PCF}$ context can separate are separated by a context that uses it. The family of closing contexts of the extended language therefore properly contains that of the base language on the shared terms, and the corresponding observational kernel is strictly finer.
\end{proof}

Finally, universality interacts with contexts in a specific way. Closure conversion and defunctionalization translate programs in a lexically scoped higher-order language into a first-order language with a single level of scope, preserving the function computed \cite{reynolds,minamide}. The nested scope tree of the source disappears as a context structure and reappears as data, in the form of explicit environment records passed as arguments.

\begin{proposition}[Contexts migrate into data; \stE]\label{prop:reify}
Closure conversion and defunctionalization are translations from lexically scoped higher-order programs to first-order programs that preserve the function computed \cite{reynolds,minamide}. In the output of either translation the source's scope tree is not represented as a context structure of the target, and the nesting survives only as the structure of data, namely explicit environment records.
\end{proposition}

The consequence for structural comparison is that the two cited translation methods do not preserve contexts as contexts, although each preserves an encoding of them; whether some other function-preserving translation could retain the context structure explicitly in its target is not excluded by the argument. Whether a common structural representation can exist even when strong equivalence fails then has a definite answer. Theorem~\ref{thm:terminalobs} supplies one for any fixed observation presheaf $\Omega$, namely the set of global sections, and it applies to universal languages as to any others. What universality does not supply is a choice of $\Omega$, and the structure retained is exactly the structure that $\Omega$ carries.
